% Source: sources/returns/S6_04_MAPS_MANIFOLDS/sections/maps_weighted_os45_invariants.tex:3–120
\paragraph{Theorem D.2.53 (Proved: complete primitive cubic and fundamental-class invariant of the weighted OS45 example).}
Take the marked OS45/default rational elliptic surface, the initial
sections \(P=2G,Q=8G\), whole-fibre weights
\((2,0,0,0,0)\) supported at \(I_8\), all multiplicative Mumford
caps, and the ordinary clutch \(e=1\). The resulting simply
connected manifold \(X\) has
\[
 H_2(X;\mathbb Z)=\mathbb Z^{15}\oplus\mathbb Z/8,\qquad
 H_3(X;\mathbb Z)=\mathbb Z/8,\qquad
 c_1(X)=c_2(X)=\lambda(X)=0.                              \tag{D.OS1}
\]
Its sixteen normalized cusp components are indexed by
\((r,s)\in\mathbb Z/8\times\mathbb Z/2\), with \(i=2r+s\).
Their divisor classes generate an index-two sublattice in
\(H^2(X;\mathbb Z)\). The full primitive lattice is generated
by \(x_0,\ldots,x_{13},y\), with
\[
             2y=\sum_{r=0}^7x_{r,1},\qquad
             \sum_{r,s}x_{r,s}=0.                        \tag{D.OS2}
\]
The complete integral cubic is (D.PC11)--(D.PC12) below.
There is a primitive cyclic coordinate \(u\in H^2(X;\mathbb Z/8)\)
with the full boundary marking (D.OF1)--(D.OF14). In the splitting
fixed by \(u\) and this primitive basis,
\[
 H_6(K(H_2(X),2);\mathbb Z)
   =\mathbb Z^{680}\oplus(\mathbb Z/8)^{120}
                       \oplus(\mathbb Z/16)^{15}
                       \oplus\mathbb Z/8,                \tag{D.OS3}
\]
and the complete fundamental-class invariant is specified by
the integral cubic and (D.OF17)--(D.OF18).
In particular
\[
 \begin{gathered}
 y^2=0,\qquad
 \langle\mathcal P_2(u)x_i,[X]\rangle=7\pmod{16}
                   \ (0\le i<14),\\
 \langle\mathcal P_2(u)y,[X]\rangle=0,\qquad
 \langle u^3,[X]\rangle=0.
 \end{gathered}                                           \tag{D.OS4}
\]
This theorem does not assert a completed smooth classification:
the refined characteristic class is still one of \(0,4\)
in the cyclic order-eight summand of \(H^4(X;\mathbb Z)\).
The full finite-coefficient invariant here does not choose
between those two possibilities.

\input{sections/fourfolds_OS45_cubic}

\paragraph{Integral vanishing of the critical class.}
The primitive-lattice argument above gives the literal relation
\(\sum x_{r,s}=0\). Therefore the target line of \(df\) is
topologically trivial. All fibres in this example are reduced
semistable fibres. The exact complex of D.2.44 consequently
applies without any additive correction: its abstract kernel
has the nowhere-zero scalar section, and its quotient line is
topologically trivial because \(c_1(X)=0\). The local unit Bott
clutch gives
\[
                     c_2(X)=\operatorname{PD}(C_{\rm crit}).
                                                               \tag{D.OS5}
\]
Here every actual complex double curve occurs once.

We retain the actual integral specialization of this complete cycle.
In the ordered bivector basis
\((e_1e_2,e_1\delta,e_1\gamma,e_2\delta,e_2\gamma,\delta\gamma)\),
the five fibre preimages, in their actual site order, are
\[
 \begin{array}{c|rrrrrr}
 I_8\text{ weighted cusp}&-16&0&6&8&0&2\\
 I_1&0&0&0&0&0&1\\
 I_1&0&8&0&0&0&1\\
 I_1&0&0&0&8&0&1\\
 I_1&0&0&0&0&0&1
 \end{array}.                                             \tag{D.OS6}
\]
For the first row, the actual local period cover has matrix
\(L=\left(\begin{smallmatrix}8&24\\-6&-16\end{smallmatrix}\right)\)
and degree sixteen. Transfer the minimal cusp's literal
three-curve preimage
\(a_1b_1+b_1a_2+a_2b_2\), expressing \(b_1,b_2\) by \(L^{-1}\)
in the actual long basis, and multiply by sixteen. This gives
the first integral row of (D.OS6), without division in homology
or an assumed covering of the global filled manifolds.
Each other row is the determinant-one transverse torus of the
corresponding primitive node, as in D.2.44.
Their sum is
\[
                        z=(-16,8,6,16,0,6).              \tag{D.OS7}
\]
In the actual twenty-three-row relation matrix \(M\) of
D.2.46, augment this vector by seventeen zeros. With columns
numbered from zero, a literal integer preimage is the column
combination with nonzero coefficients
\[
 \begin{array}{c|rrrrrrrr}
 j&6&8&10&15&16&18&19&20\\\hline
 c_j&-6&2&2&-28&2&-2&-52&-8\\
 \end{array}
 \qquad
 \begin{array}{c|rrrrrrr}
 j&26&38&41&51&56&61&66\\\hline
 c_j&2&10&10&10&10&10&10 .
 \end{array}                                             \tag{D.OS8}
\]
Direct multiplication gives \(Mc=(z,0)\).
All columns of \(M\) are the actual core prism and cap Thom
relations, so this proves that the literal critical cycle is
integrally nullhomologous. Equations (D.OS5)--(D.OS8) prove
the absolute Chern and spin-characteristic vanishing in
(D.OS1). They do not by themselves compute linking with an
odd middle-dimensional generator: local specialization chains
must be retained for that separate refined calculation.

\input{code/maps_weighted_os45_mu_proof}

% FULL OPERATIVE CONTINUATION: sources/returns/S6_04_MAPS_MANIFOLDS/sections/fourfolds_OS45_cubic.tex
% Working exact free-lattice input for the maps-owned OS45 invariant problem.
% No complete torsion fundamental-class invariant or smooth classification is asserted.
\subsubsection*{The primitive component lattice and cubic in the sixteen-component OS45 example}
We first record a saturation calculation which uses the actual component
meridians. Let $X$ be simply connected, let $N$ be a single rank-two
Mumford cusp with $D$ components, and put $E=X\setminus\operatorname{int}N$.
Assume $b_2(X)=D-1$, $H_3(N;\mathbb Z)$ is free, and
\[
 H_1(E;\mathbb Z)=\mathbb Z\oplus A,\qquad
                       \mu_i=(1,a_i),                      \tag{D.PC1}
\]
where $A$ is finite and $\mu_i$ is the positively oriented meridian of
component $i$. The component classes $x_i$ then generate a sublattice
\[
 L=\mathbb Z^D/\mathbb Z(1,\ldots,1)
                       \subset H^2(X;\mathbb Z)
\]
whose full primitive extension, inside $L\otimes\mathbb Q$, is
\[
 H^2(X;\mathbb Z)
   =L+\left\{\sum_i\widetilde{\chi(a_i)}x_i:
                         \chi\in\operatorname{Hom}(A,\mathbb Q/\mathbb Z)
        \right\}.                                        \tag{D.PC2}
\]
A tilde means any rational lift of a value in $\mathbb Q/\mathbb Z$;
changing it adds an element of $L$. In particular the component
lattice has index $|A|$.

\begin{proof}
Poincare--Lefschetz duality identifies
$H^2(X,E;\mathbb Z)=H_4(N;\mathbb Z)=\mathbb Z^D$,
with its positive component Thom basis, and
$H^3(X,E;\mathbb Z)=H_3(N;\mathbb Z)$ is free. The connecting
map from $H^1(E;\mathbb Z)$ evaluates on the meridians (D.PC1),
so its image is precisely $\mathbb Z(1,\ldots,1)$.
The relative exact sequence therefore injects $L$ into $H^2(X)$.
Since $X$ is simply connected, the latter group is free; its rank is
$D-1$ by hypothesis. Its quotient by $L$ is finite and embeds in
$H^2(E)$, whose torsion is
$\operatorname{Ext}(H_1(E),\mathbb Z)=\operatorname{Hom}(A,\mathbb Q/\mathbb Z)$.
All this torsion lifts to $H^2(X)$ because the next relative group is
free. Thus the quotient has exactly that order.

Here is also its precise embedding, rather than only its index.
For a character $\chi$ of order $m$, extend it by zero on the free
summand of $H_1(E)$. Regard it as a class in
$H^1(E;\mathbb Z/m)$ by multiplying its values by $m$.
Its relative coboundary has component coordinates
$m\widetilde{\chi(a_i)}\pmod m$. Choose the displayed integers
as an integral relative lift. Its image in $H^2(X;\mathbb Z/m)$
is zero by exactness. Hence its integral image is divisible by $m$
in $H^2(X;\mathbb Z)$, giving the class in (D.PC2).
The coefficient exact sequences identify its exterior restriction
with the corresponding torsion Bockstein, up to the harmless common
connecting-map sign. Since these characters give the full torsion
of $H^2(E)$, they give the full primitive extension. This proves
(D.PC2), including every finite denominator.
\end{proof}

Now take the exact marked data
\[
 \mathrm{OS45/default},\quad P=2G,\quad Q=8G,\quad
 (w_0,w_1,w_2,w_3,w_4)=(2,0,0,0,0),\quad e=1,              \tag{D.PC3}
\]
with all multiplicative fibres filled by their marked Mumford caps.
D.2.46 gives
\[
 H_2(X)=\mathbb Z^{15}\oplus\mathbb Z/8,\qquad
 H_3(X)=\mathbb Z/8,\qquad \pi_1(X)=1.
                                                               \tag{D.PC4}
\]
The elliptic monodromy at $I_8$ is
$\left(\begin{smallmatrix}1&8\\0&1\end{smallmatrix}\right)$
and its section cocycle is $(3,0)$. Formula (D.GI2) therefore gives
\[
 A_0=\begin{pmatrix}
 1&8&0&24\\0&1&0&0\\0&-6&1&-16\\0&0&0&1
 \end{pmatrix},\qquad
 L_{\rm per}=\left\langle(8,-6),(24,-16)\right\rangle
                 =8\mathbb Z\oplus2\mathbb Z.             \tag{D.PC5}
\]
The short basis is $(e_1,\delta)$ and the long basis is $(e_2,\gamma)$.
This is an equality of sublattices, not merely a comparison of
indices. The determinant is sixteen and the Smith factors are two
and eight. Label the normalized components by
\[
             (r,s)\in\mathbb Z/8\times\mathbb Z/2,
             \qquad i=2r+s.                              \tag{D.PC6}
\]

The actual cusp complement satisfies
\[
 H_1(E)=\mathbb Z\oplus\mathbb Z/2,
                      \qquad \mu_{r,s}=(1,s\bmod2).      \tag{D.PC7}
\]
To verify the marking, retain all peripheral generators and the
ordinary clutch in the exact degree-one presentation. Its Smith
form has seven unit factors, one factor two, and one free generator.
The torsion character in the original nine coordinates is
\[
                  (0,-2,1,2,0,0,2,-2,2)\pmod2,
                                                               \tag{D.PC8}
\]
which is precisely $\delta^*$ modulo two. The free character
evaluates to one on the cusp meridian. A height-one ray at
$(r,s)$ has literal meridian $\theta+r e_1+s\delta$; it
therefore gives exactly (D.PC7). The full original presentation,
Smith change of basis and all sixteen meridian images are saved
in the executable below.

Applying (D.PC2), there is a unique integral class $y$ with
\[
                    2y=\sum_{r=0}^{7}x_{r,1}.             \tag{D.PC9}
\]
The component lattice has exact index two. In particular it would
be incorrect to declare any fifteen of the component classes a
primitive basis. A convenient actual basis is
\[
 x_0,x_1,\ldots,x_{13},y,\qquad
 x_{14}=-2y-\sum_{\substack{i<14\\i\ \mathrm{even}}}x_i,
 \qquad
 x_{15}=2y-\sum_{\substack{i<14\\i\ \mathrm{odd}}}x_i.
                                                               \tag{D.PC10}
\]
The integral matrix of the sixteen component classes in this basis
has Smith factors $1^{14},2$ and kernel generated by the all-one
vector. These are the asserted primitive markings.

Here is the complete cubic on that basis. For component classes,
D.2.47 specializes to the following symmetric integral trilinear
form. With distinct vertices as appropriate,
\[
 \begin{aligned}
 T(x_i,x_i,x_i)&=6,\\
 T(x_i,x_i,x_j)&=
 \begin{cases}
 -2,&i=(r,s),\ j=(r,1-s),\\
 -1,&i=(r,s),\ j=(r\pm1,t),\quad t=0,1,\\
 0,&\text{otherwise},
 \end{cases}\\
 T(x_i,x_j,x_k)&=
 \begin{cases}
 1,&\{i,j,k\}\text{ is three of the four vertices in}\\[-2pt]
   &\hspace{12pt}\{(r,0),(r,1),(r+1,0),(r+1,1)\}
                   \text{ for some }r,\\
 0,&\text{otherwise}.
 \end{cases}
 \end{aligned}                                             \tag{D.PC11}
\]
All horizontal indices are cyclic modulo eight. Repeated geometric
edges and triangle cells are retained in this formula.
The remaining primitive entries are
\[
 \begin{aligned}
 T(y,x_{r,s},x_{r',s'})
 &=(-1)^s\left(
       \mathbf1_{s'=s,\ r'=r+1}
      +\mathbf1_{s'=s,\ r'=r-1}
      -2\mathbf1_{(r',s')=(r,s)}\right),\\
 T(y,y,x_i)&=0,\qquad T(y,y,y)=0.
 \end{aligned}                                             \tag{D.PC12}
\]
These formulas determine every entry on the primitive basis
(D.PC10), and every resulting entry is integral.

For a geometric verification of the half-class entries, on a
normalized $dP_6$ take the cyclic boundary classes
\[
 E_1,\ H-E_1-E_2,\ E_2,\ H-E_2-E_3,\ E_3,\ H-E_1-E_3
                                                               \tag{D.PC13}
\]
in the ray order of D.2.47. The four nonhorizontal boundary classes
sum to $2f$, where $f=H-E_1$ and $f^2=0$. Restricting (D.PC9)
therefore gives
\[
                         y|_{P_{r,s}}=(-1)^s f.            \tag{D.PC14}
\]
This division is valid in the torsion-free group $H^2(dP_6;\mathbb Z)$.
The fibre $f$ has degree one on each horizontal boundary curve
and zero on the other four. The Cartier restriction formula in
D.2.47 now gives the first line of (D.PC12). Squaring (D.PC14)
gives $y^2x_i=0$; multiplying (D.PC9) by $y^2$ gives $y^3=0$.
Thus no rational cubic has been silently substituted for the
integral form.

Finally $c_1(X)=0$ and both $c_2$ and $p_1$ pair to zero with
every class in the primitive free basis. For $c_1$, the normalized
component tangent calculation is zero and the component restriction
is rationally injective by the connected toric graph Laplacian;
$H^2(X)$ is torsion-free, so this proves integral vanishing.
For $c_2,p_1$, D.2.47 gives zero on every component. The component
classes have finite index, and the evaluations are integers, so they
also vanish on $y$. These statements concern the free characteristic
pairings; they do not determine the torsion class of $c_2$, the
refined class, or the full finite-coefficient fundamental-class
invariant.

\path{code/fourfolds_OS45_primitive_cubic.py}
reconstructs the actual complement before using its torsion
character, checks all sixteen meridians, computes the exact
index-two primitive lattice, and verifies all $15^3=3375$ primitive
triples and all $16^3=4096$ recovered component triples. The
normalization restriction (D.PC14), the actual period basis, and
all sparse integral entries are retained in its JSON output.
The preliminary component-only calculation is retained separately
in \path{code/fourfolds_OS45_component_ring.py}.

% FULL OPERATIVE CONTINUATION: sources/returns/S6_04_MAPS_MANIFOLDS/code/maps_weighted_os45_mu_proof.tex
% D.2.53: independently audited complete fundamental-class calculation.
% The refined class is explicitly a separate remaining calculation.
\subsubsection*{The actual torsion coordinate and complete fundamental-class
invariant of the sixteen-component OS45 manifold}
We use exactly the datum (D.PC3), its period marking (D.PC5), and the
primitive free basis (D.PC10). Put \(k=\mathbb Z/8\).
There is a primitive torsion coordinate
\(u\in H^2(X;k)\) whose complete cusp support is
\[
 \Psi(r,s)=4\left\lfloor r/8\right\rfloor
       +3(r\bmod8)\left\lfloor s/2\right\rfloor
                                      \quad\hbox{in }k.       \tag{D.OF1}
\]
Here and below \(r\bmod8\) means the representative between zero and
seven. The assertion includes the entire boundary cohomology class,
including its coefficient-extension part. We prove this before using
any support-function evaluations.

\paragraph{The specified core cochain.}
Use the actual degree-two core model of D.2.46. Its fibre cells are
ordered
\[
 e_1e_2,\ e_1\delta,\ e_1\gamma,\ e_2\delta,\
                                      e_2\gamma,\ \delta\gamma.
\]
For the six boundary monodromies, including the ordinary identity cap,
let
\[
 D=(A_0-I,\ldots,A_4-I),\qquad K=\ker_{\mathbb Z}D.
\]
Choose the integral cycle basis and the integral left inverse \(K^+\)
retained by the exact program. The ordinary degree-two cycles are the
six fibre cells followed by this graph-cycle basis. The actual
relation matrix \(M\) has twenty-three rows. If
\(UMV\) is its stored Smith form, the unique nonunit factor is eight.
Let \(\varphi\) be the negative of the corresponding row of \(U\).
Then
\[
 \varphi M=0\pmod8,\qquad
 \varphi|_{\Lambda^2\Lambda}=(1,0,0,0,0,0).                \tag{D.OF2}
\]
Its value on the stored Smith torsion generator is \(-1\).
This is already a primitive functional on the order-eight image of
the core in \(H_2(X)\).

We specify the cochain, rather than merely its homology functional.
The inclusion \(K\subset\mathbb Z^{20}\) is primitive because it is
the kernel of an integral map. Extend the graph row by
\(\varphi_{\rm graph}K^+\), and use the fibre row in (D.OF2).
This gives a literal degree-two cochain on the core mapping-cylinder
model. The degree-three boundary columns are the actual exterior
monodromy relations, so (D.OF2) says that it is closed modulo eight.
The fixed bar-prism comparisons in D.2.46 identify this cochain with
one on the actual torus bundle. In particular its restriction is not
left unspecified by the possible \(2\)-torsion in \(H_1\) of the core.

At the rank-two cusp the actual peripheral inclusion has graph
matrix \(J=K^+J_{\rm raw}\). The retained cochain gives
\[
 \varphi_{\rm graph}J=(0,52,0,148)
       \quad\hbox{on }(e_1,e_2,\delta,\gamma).             \tag{D.OF3}
\]
The quadratic correction on each of these four basis loops is zero
at this cusp. Thus the two long-strip constants are
\[
                      (c_1,c_2)=(4,4)\quad\hbox{in }k.   \tag{D.OF4}
\]
The full Smith row, the chosen left inverse, and the actual
peripheral matrices are reproducible by
\path{code/maps_weighted_os45.py}. In particular (D.OF3) is not
obtained by restricting only the regular-fibre alternating form.

\paragraph{A full boundary comparison for finite-coefficient supports.}
We give the comparison explicitly, since this cusp has torsion in
its boundary \(H_1\). Let \(\widetilde N\) be the common universal
cover of the triangular Mumford cusp and let
\(\widetilde M\) be its boundary. Its normalized components are
indexed by \(v\in\mathbb Z^2\). The boundary cover retracts to the
three-torus generated by the two short angular loops and the positive
\(q\)-meridian \(\theta\). D.2.43 gives
\[
 H_1(\widetilde N;\mathbb Z)=0,\quad
 H_2(\widetilde N;\mathbb Z)=I^2,\quad
 \operatorname{im}\{H_2(\widetilde M;\mathbb Z)\to
                         H_2(\widetilde N;\mathbb Z)\}=0. \tag{D.OF5}
\]
Indeed the three boundary generators are the
angular area and the two angular-normal tori, all with the actual
bounding sweeps described there.

The relative component Thom cochains identify
\[
 H^2(\widetilde N,\widetilde M;k)
                  =\operatorname{Map}(\mathbb Z^2,k).
\]
One can obtain this directly from locally finite
component chains: a degree-four locally finite cycle is an arbitrary
family of fundamental cycles of the normalized components, and
degree-four intersection strata do not occur. Equivalently,
Poincare--Lefschetz duality uses Borel--Moore degree-four homology.
No compact-support convention replaces the displayed product of
coefficient groups by a direct sum.

The positive meridian of component \(v\) is
\(\theta+v_1e_1+v_2\delta\). Therefore the relative connecting map
from \(H^1(\widetilde M;k)\) sends a covector \((c,m_1,m_2)\) to
the affine function \(c+m_1v_1+m_2v_2\). By (D.OF5), the relative
sequence is exactly
\[
 0\longrightarrow\operatorname{Aff}(\mathbb Z^2,k)
 \longrightarrow\operatorname{Map}(\mathbb Z^2,k)
 \longrightarrow H^2(\widetilde N;k)\longrightarrow0.     \tag{D.OF6}
\]
Here \(\operatorname{Aff}\) consists of these three-parameter
affine functions. This also identifies the quotient with
\(\operatorname{Hom}(I^2,k)\), consistently with the actual
normalization computation of D.2.43.

Let the actual period lattice \(L\) act by translation on the
argument \(v\). If a function \(\Psi\) satisfies
\[
                  \Psi(v+l)-\Psi(v)=a_l(v)
                       \in\operatorname{Aff},           \tag{D.OF7}
\]
its absolute class on \(\widetilde N\) is \(L\)-invariant.
It descends to \(N_L\): in total degree two the covering
cochain double complex has \(H^1(\widetilde N;k)=0\), and
the only possible obstruction lies in \(H^3(L;k)=0\).
The choices of descent form a torsor for \(H^2(L;k)=k\).

There is a useful precise version of the boundary part of this
argument. The boundary group is the split extension
\[
       \pi_1(M_L)=\mathbb Z^3\rtimes L,                    \tag{D.OF8}
\]
where the normal \(\mathbb Z^3\) is generated by the two shorts
and \(\theta\), and the two long loops give the specified splitting.
Its dual module \(H^1(\widetilde M;k)\) is exactly
\(\operatorname{Aff}\): evaluating a covector on
\(\theta+v\) produces its affine function, and translating \(v\)
gives the deck action.
For boundary classes with zero restriction to \(\widetilde M\),
the degree-two filtration is the short exact sequence
\[
 0\longrightarrow H^2(L;k)\longrightarrow
 \ker\{H^2(M_L;k)\to H^2(\widetilde M;k)\}
 \longrightarrow H^1(L;\operatorname{Aff})\longrightarrow0.
                                                               \tag{D.OF9}
\]
The splitting of (D.OF8) kills the degree-two transgression;
there is no degree-three group cohomology of \(L\).
These two observations also prove (D.OF9) directly from the
two-dimensional free resolution of \(L=\mathbb Z^2\).
Restriction to the actual long two-torus splits its leftmost term.

The image in the rightmost group of the boundary value of a
descended support is represented by the full affine cocycle
\(l\mapsto a_l\) in (D.OF7). This follows either by chasing
(D.OF6) in that resolution, or by choosing the relative Thom
cochain \(\Psi\): its deck difference is the boundary one-cocycle
whose value on \(\theta+v\) is precisely \(a_l(v)\).
Thus the gradients, constants, and one long-torus value determine
the entire boundary class. This statement retains every
\(\operatorname{Ext}(H_1(M_L),k)\) ambiguity.

In our positive fan convention a core cochain with short-long
matrix \(E\), zero angular area and zero short-normal strip values
has affine cocycle
\[
 a_{l_j}(v)=-E(v,b_j)+c_j,\qquad j=1,2,                  \tag{D.OF10}
\]
where \(b_1=e_2,b_2=\gamma\) and \(l_j=(A-I)b_j\).
For a direct verification, translate the marked \(\theta,b_j\)
strip from component zero to component \(v\). Its change is the
oriented short-long torus with value \(-E(v,b_j)\).
The initial strip contributes \(c_j\). There is no additional
angular-area term, since that coefficient is zero. The literal
basis-loop prism correction is the one already included in
(D.OF3). This gives (D.OF10) at the cochain level, including
its constant term.

For (D.OF2)--(D.OF4), \(E=\left(\begin{smallmatrix}1&0\\0&0
\end{smallmatrix}\right)\). The required affine functions are
therefore
\[
             a_{(8,-6)}=-r+4,\qquad a_{(24,-16)}=4.
                                                               \tag{D.OF11}
\]
The support (D.OF1) has exactly these differences.
This is an identity for all integers, not a bounded numerical
test: write \(r=8a+r_0,s=2b+s_0\) with \(0\le r_0<8\) and
\(0\le s_0<2\), and substitute. The two differences are
\(4-9r_0\) and \(12-24r_0\), respectively, which reduce to
(D.OF11).
Choose the remaining descent parameter so its value on the
actual long torus \(\langle e_2,\gamma\rangle\) is zero,
as it is for (D.OF2).
Equations (D.OF9)--(D.OF11) now prove equality on the full
cusp boundary, including both its finite-coefficient and
integral-torsion contributions.

For an independent check of the short-long marking, retain the
actual annulus map from D.2.43, not only its augmentation.
If \(x^py^q-1=(x-1)S_x+(y-1)S_y\), the angular vector \((a,b)\)
closes to the \(I^2\)-class with coefficients
\[
             (aS_x,\ bS_x+aS_y,\ bS_y)                   \tag{D.OF12}
\]
on \(((x-1)^2,(x-1)(y-1),(y-1)^2)\).
Its literal Picard columns are
\[
 x^{-1}y^{-2}E_2+y^{-1}E_3,\quad
       -x^{-1}y^{-1}E_2,\quad
       x^{-1}E_1+x^{-2}y^{-1}E_2.                        \tag{D.OF13}
\]
In the present positive component labels \(x\) translates by
\((-1,0)\), and \(y\) by \((0,1)\); the path and angular exponents
are changed simultaneously. Substituting (D.OF1) into all the
shifted columns gives the integer values
\[
                         \begin{pmatrix}9&24\\0&0\end{pmatrix}
                  \equiv E\pmod8.                       \tag{D.OF14}
\]
The path dependence in \(S_x,S_y\) is exactly an
Hilbert--Burch relation, so this checks the actual annulus
marking rather than a rational replacement.

\paragraph{Extension through the other caps.}
Each of the four other multiplicative fibres has primitive
rank-one monodromy. Its boundary \(H_1\) is free; so is the
ordinary cap boundary \(H_1\).
The regular-fibre specialization is onto the degree-two group
of each such cap: in a primitive nodal marking this is the
specialization from \(T^2\times T^2\) to
\((S^2\vee S^1)\times T^2\); for the ordinary cap it is the
identity on the central \(T^4\).
Their actual degree-two boundary kernels
are precisely the columns used in \(M\), including the
ordinary cap's affine-meridian column.
Equation (D.OF2) therefore makes the specified core cochain
extend across each such cap. The free boundary \(H_1\)
ensures that agreement of the resulting integral-cycle
functional is agreement of the full finite-coefficient
cohomology class. These choices and the cusp comparison
glue to the asserted \(u\in H^2(X;k)\).
Its torsion value is the unit \(-1\), unchanged from (D.OF2);
we choose the opposite torsion generator to normalize it to one.
Choose free lifts of the homology basis dual to (D.PC10)
which are annihilated by \(u\). This fixes the splitting
\[
                         H_2(X)=\mathbb Z^{15}\oplus\mathbb Z/8
                                                               \tag{D.OF15}
\]
used for the fundamental-class coordinates.

\paragraph{The complete coordinate values.}
Let \(R_0,\ldots,R_5\) be the six positive fan rays in cyclic
order and define the actual wall values
\[
 d_j(v)=\Psi(v)+\Psi(v+R_j)
                    -\Psi(v+R_{j-1})-\Psi(v+R_{j+1}).
\]
An integral representative for \(u|_{P_v}\) is
\[
 \eta_v=(d_0+d_1+d_2)H-d_0E_1-d_2E_2-d_4E_3.             \tag{D.OF16}
\]
Opposite wall restrictions agree modulo eight. Changes of integral
representative by eight times a class preserve its square modulo
sixteen. The full exact component table is retained in
\path{records/maps_weighted_os45_mu.json}.

For component labels \(i,j\), all mixed ordinary coordinates are
specified without a long matrix display by
\[
 \begin{aligned}
 \langle ux_i^2,[X]\rangle
                    &=-\sum_{a=0}^{5}d_a(v_i)\pmod8,\\
 \langle ux_ix_j,[X]\rangle
                    &=\sum_{\{a:\ v_i+R_a=v_j\ {\rm mod}\ L\}}
                                       d_a(v_i)\pmod8\quad(i\ne j),\\
 \langle ux_i y,[X]\rangle
                    &=(-1)^{s_i}\eta_i\cdot(H-E_1)\pmod8.
 \end{aligned}                                             \tag{D.OF17}
\]
The last value is four for component labels \(0,1,14,15\)
and zero for the other twelve. These formulas use actual
Cartier restrictions and retain the repeated vertical edges.
On the primitive free basis \(x_0,\ldots,x_{13},y\), the remaining
power coordinates are
\[
 \begin{split}
 \langle\mathcal P_2(u)x_i,[X]\rangle&=7\pmod{16}
                                             \quad(0\le i<14),\\
 \langle\mathcal P_2(u)y,[X]\rangle&=0\pmod{16},\\
 \langle u^3,[X]\rangle&=0\pmod8,\qquad
 \langle uy^2,[X]\rangle=0\pmod8.
 \end{split}                                               \tag{D.OF18}
\]
The first line is \(\eta_i^2\) on the normalized component.
The two omitted component labels have square fifteen modulo
sixteen, consistently with their primitive-basis relations.
The remaining three evaluations, including the full
order-sixteen value which cannot be obtained by division,
are proved in the following localization input.
\input{sections/fourfolds_OS45_localization}

In particular \(y^2=0\) integrally. Its free part is zero by
(D.PC12), and its torsion part is detected injectively by the
primitive \(u\); the last equality of (D.OF18) kills that part.

Finally D.2.48's primary coefficient-group calculation and
the integral Kunneth decomposition give
\[
 H_6(K(\mathbb Z^{15}\oplus\mathbb Z/8,2);\mathbb Z)
       =\mathbb Z^{680}\oplus(\mathbb Z/8)^{120}
                         \oplus(\mathbb Z/16)^{15}
                         \oplus\mathbb Z/8.              \tag{D.OF19}
\]
The free coordinates are exactly the primitive integral cubic
(D.PC11)--(D.PC12). Equations (D.OF17)--(D.OF18) give all
remaining coordinates: the first mixed group is detected by
\(ux_ix_j\), the second by the full Pontryagin square, and the
last by \(u^3\). Hence this determines the complete
fundamental-class invariant \(\mu\), not merely its image
under ordinary triple cup products.
The refined class remains a separate \(0\) versus \(4\)
calculation until the actual odd middle-dimensional generator
and its reference framing have been evaluated.

\path{code/maps_weighted_os45_mu.py} checks the exact Smith
cochain, both full affine strip constants, all sixteen wall
restrictions, the full shifted annulus map, every primitive
mixed entry and the complete moduli in (D.OF18).
The independent fourfolds programs check all primitive cubic
entries and all-character localization identities.

% FULL OPERATIVE CONTINUATION: sources/returns/S6_04_MAPS_MANIFOLDS/sections/fourfolds_OS45_localization.tex
% Auxiliary input to the maps-owned OS45 fundamental-class calculation.
% The full boundary comparison identifying the support with global u is required.
\subsubsection*{The remaining finite-coefficient evaluations for OS45}
Let $u\in H^2(X;\mathbb Z/8)$ be the actual primitive torsion
coordinate constructed by the complete boundary comparison, with
cusp support
\[
 \Psi(r,s)=4\left\lfloor\frac r8\right\rfloor
       +3(r\bmod8)\left\lfloor\frac s2\right\rfloor
                           \quad\pmod8.                  \tag{D.OL1}
\]
Here $r\bmod8$ is the representative from zero to seven. This
subsection uses that comparison on the whole cusp boundary;
matching only the regular-fibre bilinear form would not suffice.
The actual affine changes are
\[
 \Delta_{(8,-6)}\Psi=-r+4,\qquad
 \Delta_{(24,-16)}\Psi=4\quad\pmod8.                     \tag{D.OL2}
\]
In particular both have zero slope in the scalar $s$ direction.
Let $y$ be the primitive class in (D.PC9).
We prove the three full evaluations
\[
 \langle u^3,[X]\rangle=0\pmod8,\qquad
 \langle\mathcal P_2(u)y,[X]\rangle=0\pmod{16},\qquad
 \langle uy^2,[X]\rangle=0\pmod8.                         \tag{D.OL3}
\]
The second formula includes the order-two part which cannot be
recovered by halving the sum of the component evaluations.

The scalar circle action is semifree. Its fixed locus consists of
the sixteen vertical toric curves $C_{r,s}$ joining vertices
$(r,s)$ and $(r,s+1)$, with $(r,s)$ indexed as in (D.PC6).
Each has normal bundle $\mathcal O(-1)\oplus\mathcal O(-1)$
with weights $1,-1$. In every other triangular toric chart the
nonzero weights are also primitive. At each of the four zero-free
$I_1$ caps the primitive angular collapsing vector has first
coordinate one in its elliptic/scalar plane; the scalar circle
therefore projects primitively to the transverse torus and is
free. The ordinary torus cap is free as well. Thus there are no
unrecorded finite-stabilizer strata in this localization.

Since $X$ is simply connected, both $u$ and $y$ have circle-equivariant
lifts. Indeed the only outgoing differentials from ordinary degree
two in the Borel spectral sequence have targets involving
$H^1(X)$ or $H^3(BS^1)$, both zero. Let $x$ be the positive
generator of $H^2(C_{r,s};\mathbb Z)$ and let $t$ be the generator
of $H^2(BS^1;\mathbb Z)$. Write
\[
 \widehat u|_{C_{r,s}}=d_{r,s}x+\lambda_{r,s}t,
 \qquad
 \widehat y|_{C_{r,s}}=\mu_{r,s}t.                         \tag{D.OL4}
\]
There is no $x$ term in the second expression: (D.PC14) and
$f\cdot C_{r,s}=0$ give $y(C_{r,s})=0$.

For clarity the ordinary degrees and the fibre characters are
computed separately. The degree of the support on a vertical wall is
\[
 d_{r,s}=\Psi(r,s)+\Psi(r,s+1)
              -\Psi(r+1,s)-\Psi(r-1,s+1)\pmod8 .         \tag{D.OL5}
\]
Choosing representatives from zero to seven gives
\[
\begin{array}{c|cc}
 r&d_{r,0}&d_{r,1}\\\hline
 0,7&4&7\\
 1,\ldots,6&0&3.
\end{array}                                                \tag{D.OL6}
\]
The scalar fibre character supplied by the support is its vertical
slope, up to a single common sign. Thus one may take
\[
                 \lambda_{r,0}=0,\qquad
                 \lambda_{r,1}=3r .                      \tag{D.OL7}
\]
This is periodic under the actual period lattice because of
(D.OL2). To justify the common character, rather than making a
separate choice on each normalized component, compare along the
actual toric orbit curves. Their degrees determine the differences
of the characters at their fixed endpoints. Horizontal and diagonal
orbit curves connect the sixteen fixed curves in a connected graph:
they join $(r,s)$ to $(r+1,s-1)$ and to $(r-1,s)$ respectively.
The support realizes all these differences. Consequently any global
equivariant lift has the characters in (D.OL7), with a common sign
convention and one common additive character $L\in\mathbb Z/8$.
No global integral line lifting $u$ is assumed.

The integral component divisors have their canonical equivariant
linearizations. Along $C_{r,s}$ the two adjacent component normal
weights are $-1,+1$ in the lower/upper vertex order. Thus the weight
of the odd-component sum in (D.PC9) is $+1$ for $s=0$ and $-1$
for $s=1$. The difference
$2\widehat y-\sum_r\widehat x_{r,1}$ is a global multiple $ct$.
The existence of integral fixed-point weights forces $c$ odd;
changing the lift of $y$ changes $c$ by an even integer. Normalize
$c=-1$. Then
\[
                       \mu_{r,0}=0,\qquad
                       \mu_{r,1}=-1.                     \tag{D.OL8}
\]
Every other lift adds one common integer $V$.

Each restriction of $\widehat u$ in (D.OL4) lifts integrally on
$C_{r,s}\times BS^1$. Thomas's cochain formula therefore gives
$\mathcal P_2(\widehat u)=(d x+\lambda t)^2$ with the full
modulus sixteen \cite[equations (2.1), (3.4), pp.267--268]{ffThomas56}.
The ordinary cube uses modulus eight. The equivariant normal Euler
class is
\[
                           (-x+t)(-x-t)=-t^2.
\]
Localization is valid over these coefficient rings because the
complement of the fixed locus is free and all normal weights are
units. Taking the coefficient of $x$ after division by $-t^2$ gives
\[
 \begin{aligned}
 \langle u^3,[X]\rangle&=-3\sum_{r,s}d_{r,s}\lambda_{r,s}^2
                                                        &&\pmod8,\\
 \langle\mathcal P_2(u)y,[X]\rangle
       &=-2\sum_{r,s}d_{r,s}\lambda_{r,s}\mu_{r,s}
                                                        &&\pmod{16},\\
 \langle uy^2,[X]\rangle&=-\sum_{r,s}d_{r,s}\mu_{r,s}^2
                                                        &&\pmod8.
 \end{aligned}                                             \tag{D.OL9}
\]
The degree-zero base ring injects into its localization, so these
are identities before inverting $t$ as well.

The representatives (D.OL6)--(D.OL8) satisfy the exact integer sums
\[
 \sum d_{r,s}=40,\qquad \sum d_{r,s}\lambda_{r,s}=336,
 \qquad \sum d_{r,s}\lambda_{r,s}^2=5544,
 \qquad \sum_r d_{r,1}=32.                               \tag{D.OL10}
\]
The three values in (D.OL9) are consequently $-16632$, $672$,
and $-32$, proving (D.OL3) with their full moduli.
For completeness, under the replacements
$\lambda\mapsto L+\epsilon\lambda$ and $\mu\mapsto V+\mu$,
where $\epsilon=\pm1$, the exact values are
\[
 \begin{aligned}
 -120L^2-2016\epsilon L-16632,\\
 -80LV+64L-672\epsilon V+672\epsilon,\\
 -40V^2+64V-32.
 \end{aligned}                                             \tag{D.OL11}
\]
Every coefficient is divisible by eight, sixteen, and eight,
respectively. Thus the results are independent of all global
linearization choices, not only of the displayed normalization.

The same support gives all component mixed evaluations. For a
vertex $v$, let $R_j$ be the six cyclic rays of D.2.47 and put
\[
 d_j(v)=\Psi(v)+\Psi(v+R_j)
                  -\Psi(v+R_{j-1})-\Psi(v+R_{j+1}).
\]
An integral representative of its component restriction is
\[
 \eta_v=(d_0+d_1+d_2)H-d_0E_1-d_2E_2-d_4E_3.
                                                               \tag{D.OL12}
\]
Changing the representative by eight times an integral class
changes neither its square modulo sixteen nor any of its needed
pairings modulo eight. In the primitive basis of (D.PC10),
\[
 \langle\mathcal P_2(u)x_i,[X]\rangle=7\pmod{16}
                     \quad(0\le i<14),
\]
while the two omitted components $i=14,15$ have value fifteen.
Also $\langle ux_i y,[X]\rangle$ equals four for $i=0,1,14,15$
and zero for all the other component labels. The remaining
$\langle ux_i x_j,[X]\rangle$ are exactly
$\eta_i\cdot\nu_i^*D_j\pmod8$, with the explicit Cartier
restriction in D.2.47. Together with (D.OL3) these determine the
finite-coordinate evaluations once the support comparison is
established.

The exact program
\path{code/fourfolds_OS45_localization.py}
retains every wall value and every $\eta_v$, verifies the actual
quasiperiods on all sixteen residues, and proves the six
all-character polynomial divisibility identities (D.OL11).
It explicitly records the required geometric support comparison
as a separate hypothesis.
